\section{Tinjauan Pencapaian Kompetensi Secara Menyeluruh}

\subsection{Pemetaan Asesmen ke CPMK}

\begin{table}[htbp]
\centering
\begin{tabular}{|>{\raggedright\arraybackslash}p{5cm}|>{\raggedright\arraybackslash}p{8cm}|}
\hline
\textbf{CPMK / Sub-CPMK} & \textbf{Diukur Melalui} \\
\hline
CPMK-1: Sub-CPMK 1.1--1.3 (Proposisi, operator, tabel kebenaran) & Pilihan Ganda (1--2), Bagian C Tabel Kebenaran \\
\hline
CPMK-2: Sub-CPMK 2.1--2.3 (Translasi, ekuivalensi, tautologi) & Pilihan Ganda (3), Essay (1), Bagian C (1) \\
\hline
CPMK-3: Sub-CPMK 3.1, 3.2 (Logika predikat, inferensi) & Pilihan Ganda (4--5), Essay (2, 5) \\
\hline
CPMK-4: Sub-CPMK 4.1, 4.2, 4.3 (Metode bukti, induksi, Prolog) & Pilihan Ganda (6--7), Essay (3--5), Bagian C (2--3), Bagian D \\
\hline
\end{tabular}
\caption{Pemetaan Komponen Asesmen ke CPMK dan Sub-CPMK}
\end{table}

\subsection{Self-Assessment Checklist}

Sebelum mengerjakan asesmen akhir, pastikan Anda telah menguasai:

\begin{checklist}
  \item Proposisi dan operator logika (negasi, konjungsi, disjungsi, implikasi, bikondisional)
  \item Tabel kebenaran dan ekuivalensi logika
  \item Tautologi, kontradiksi, dan kontingensi
  \item Translasi bahasa alami ke simbol logika proposisional dan predikat
  \item Kuantifier universal dan eksistensial beserta negasinya
  \item Modus ponens, modus tollens, silogisme, dan aturan inferensi lain
  \item Bukti langsung, kontraposisi, dan kontradiksi
  \item Induksi matematika (basis dan langkah induktif)
  \item Dasar Prolog: fakta, aturan, query, dan unifikasi
\end{checklist}
